\documentclass[a4paper,10pt]{amsart}
% Complete shared setup. Load after \documentclass.
\usepackage{iftex}
\RequireLuaTeX

% myMath loads mathtools before unicode-math, as required by unicode-math.
\usepackage{myMath}
\usepackage{fontspec}
\usepackage{unicode-math}

% Use installed STIX Two fonts by family name.
\setmainfont{STIX Two Text}
\setmathfont{STIX Two Math}[StylisticSet={4,10,11}]

% Page layout and general-purpose packages.
\usepackage[headings]{fullpage}
\usepackage{xcolor}
\usepackage[shortlabels]{enumitem}
\usepackage{float}
\usepackage{algorithm}
\usepackage[noend]{algpseudocode}
\usepackage{multicol}
\usepackage{xfrac}
\usepackage{microtype}
\setlength{\marginparwidth}{2cm}
\relpenalty=9999
\binoppenalty=9999
\setlist[enumerate,1]{label=(\roman*),font=\upshape}

% Diagrams.
\usepackage{tikz}
\usepackage{pgfplots}
\pgfplotsset{compat=1.18}
\usetikzlibrary{arrows.meta,positioning}
\tikzset{
  every path/.style={>={Stealth[length=4pt,width=5pt,inset=2.3pt,sep=-1.7pt]}}
}

% Revision colours.
\definecolor{fresh}{HTML}{2BB101}
\definecolor{checked}{HTML}{1E5E06}
\definecolor{double}{HTML}{5E3800}
\definecolor{external}{HTML}{A81A78}
\definecolor{later}{HTML}{0410FF}
\definecolor{minor-rev}{HTML}{D96A09}
\definecolor{major-rev}{HTML}{C90000}
\definecolor{skip}{HTML}{FFFFFF}
\definecolor{normal}{HTML}{000000}

% Bibliography resource and style.
\usepackage[
  backend=biber,
  maxbibnames=99,
  datamodel=mrnumber,
  giveninits=true,
  style=alphabetic,
  sortcites,
  sorting=nyt
]{biblatex}
\addbibresource{refs.bib}
\renewcommand*{\multicitedelim}{\addcomma\space}
\renewcommand*{\bibfont}{\footnotesize}
\DeclareFieldFormat[article,inproceedings,incollection]{title}{\mkbibemph{#1}}
\DeclareFieldFormat{journaltitle}{#1}
\DeclareFieldFormat[article]{volume}{\mkbibbold{#1}}
\renewbibmacro{in:}{%
  \ifentrytype{article}{}{\printtext{\bibstring{in}\intitlepunct}}}
\DeclareFieldFormat{pages}{#1}
\DeclareFieldFormat{mrnumber}{%
  MR\addcolon\space
  \ifhyperref
    {\href{https://www.ams.org/mathscinet-getitem?mr=#1}{\nolinkurl{#1}}}
    {\nolinkurl{#1}}}
\renewbibmacro*{doi+eprint+url}{%
  \iftoggle{bbx:doi}{\printfield{doi}}{}%
  \newunit\newblock
  \printfield{mrnumber}%
  \newunit\newblock
  \iftoggle{bbx:eprint}{\usebibmacro{eprint}}{}%
  \newunit\newblock
  \iftoggle{bbx:url}{\usebibmacro{url+urldate}}{}}
\DeclareSourcemap{
  \maps[datatype=bibtex]{
    \map{\step[fieldset=issn,null]}
  }
}
\AtEveryBibitem{\iffieldundef{doi}{}{\clearfield{url}}}
\AtEveryCitekey{\iffieldundef{doi}{}{\clearfield{url}}}

% Footnotes and links. Load cleveref before declaring theorem environments.
\usepackage[perpage]{footmisc}
\renewcommand{\thefootnote}{%
  \ifcase\value{footnote}%
  \or\textdagger
  \or\textdaggerdbl
  \fi
}
\usepackage{hyperref}
\hypersetup{
  hypertexnames=false,
  colorlinks=true,
  linkcolor=blue,
  citecolor=green,
  filecolor=magenta,
  urlcolor=cyan
}
\usepackage[capitalise,noabbrev]{cleveref}
\usepackage{myEnvs}
% amsart's default proof marker otherwise uses a legacy AMS symbol font.
\renewcommand{\qedsymbol}{\ensuremath{\mdlgwhtsquare}}

% Additional notation shared with the Jung paper.
\DeclareMathOperator{\deff}{def}
\renewcommand{\vec}[2]{%
  \begin{pmatrix}
    #1\\
    #2
  \end{pmatrix}%
}


\begin{document}




\title{A simple algebraic proof of Jung's theorem}

%    Only \author and \address are required; other information is
%    optional.  Remove any unused author tags.
\author{Kacper Proniewski}
\address{Institute of Mathematics,
  Jagiellonian University, Kraków, Poland}
\email{kacper.proniewski@student.uj.edu.pl}

\date{\today}

\begin{abstract}
  We give a short and elementary proof of Jung's theorem, which states that if \(K\) is a field of characteristic 0, then the group of automorphisms \(\Aut_K K[x,y]\) is generated by
  linear automorphisms
  and elementary automorphisms. The core argument is based on \autocite{ML2014}. Here, we present it in a more self-contained, accessible, and streamlined manner.
\end{abstract}

\maketitle

%%%%%%%%%%%%%%%%%%%%%%%%%%%%%%%%%%%%%%%%%%%%%%%%%%%%%%%%%%%%%%%%%%%%%%%%%%%%%%%%%%%%%%%%%%%%%%%%%%%%%%%%%%%
\section{Introduction}
Let \(K\) be a field of characteristic 0. In 1942 Jung
\autocite{J42} showed that if \(K\) is algebraically closed, then any automorphism of \(R\coloneqq K[x,y]\)
is a composition of elementary automorphisms
(given by \(x\mapsto x\), \(y\mapsto y+p(x)\) or \(x\mapsto x+p(y)\), \(y\mapsto y\) for a polynomial \(p\))
and linear automorphisms. Later van der Kulk proved this theorem for an arbitrary field \autocite{vdK53}.
Since then, many authors have given proofs of this result
\autocite{SL02, AM75, Gut61, MW88, Nag72, NVC03, Ren68, VGG14}.
The approach we present here is a streamlined version of the argument given in the note by Makar-Limanov \autocite{ML2014}, which studies more generally the possible counterexamples to the
Jacobian conjecture.
The theorem is true for an arbitrary field, even of positive characteristic. For a proof in this case, see for example the geometric argument presented in \autocite{SL02}.

The key step in our proof is to find a polynomial \(\zeta=ax(y+bx^c)\) such that \(J(\zeta, f_w)=f_w\), where \(J(f,g)\) stands for the Jacobian determinant of two polynomials, and \(f_w\) is 
the leading form of \(f\) with respect to the weight \(w\).
The argument is completely elementary, using only basic facts about Newton polygons, derivations, and homogeneous polynomials. In \cref{sec:2} 
we prove these preliminary results, and in \cref{sec:3} we use them to prove Jung's theorem.

%%%%%%%%%%%%%%%%%%%%%%%%%%%%%%%%%%%%%%%%%%%%%%%%%%%%%%%%%%%%%%%%%%%%%%%%%%%%%%%%%%%%%%%%%%%%%%%%%%%%%%%%%%%
\section{Preliminaries}\label{sec:2}
In this section we introduce the basic tools needed in the proof. We define weighted degrees of polynomials, Newton polygons, and present some basic results in the theory of derivations.
\begin{notation}
  We let \(J(f,g)\) denote the Jacobian determinant of two polynomials \(f,g\in K[x,y]\).
  Unless stated otherwise, \(K\) is a field of characteristic 0 and \(R=K[x,y]\).
\end{notation}
\begin{definition}[Weight function]
  Given \(\a,\b\in \Z\) we define a \emph{weight function} \(w\colon R\setminus\set{0}\to \Z\) on monomials by \(w(x^iy^j)=\a i+\b j\) and on a general nonzero polynomial \(p=\sum_{i,j}p_{ij}x^iy^j\) by
  \[w(p)=\max\set{w(x^iy^j)\mid p_{ij}\ne 0}.\]
  We call \(w(p)\) the \(w\)-degree of \(p\).
\end{definition}
Each polynomial \(p\) can be written as \(\sum_i p_i\), where \(p_i\) are \(w\)-homogeneous polynomials (by grouping monomials with the same \(w\)-degree).
We let \(p_w\) denote the leading form of the polynomial \(p\) with respect to the weight \(w\).
\begin{proposition}\label{prop:jacob-w}
  Let \(w\) be a weight function on \(R\) and let \(f,g\in R\) be polynomials such that \(J(f,g)\ne 0\). Then
  \[w(J(f,g))\le w(f)+w(g)-w(x)-w(y).\]
  Moreover, if \(f,g\) are \(w\)-homogeneous, then equality holds.
\end{proposition}
\begin{proof}
  Taking a partial derivative with respect to \(x\) (resp. \(y\)) lowers the \(w\)-degree by \(w(x)\) (resp. \(w(y)\)). Upon computing the Jacobian both \(f_xg_y\) and \(f_yg_x\) terms 
  have the same weight \(w(f)+w(g)-w(x)-w(y)\), this proves the first statement, since \(w(p+q)\le \max\set{w(p),w(q)}\). The second statement follows, since by homogeneity either
  all the monomials in the expression \(J(f,g)=f_xg_y-f_yg_x\) will cancel, or the degree of \(J(f,g)\) will be equal to \(w(f)+w(g)-w(x)-w(y)\).
\end{proof}
\begin{definition}[Newton polygon]
  For any \(f\in R\) of the form \(f=\sum_{i,j}f_{ij}x^iy^j\) we define
  \[\supp f=\set{(i,j)\in \Z^2\mid f_{ij}\ne 0}.\]
  Then the \emph{Newton polygon} \(\cN(f)\) is defined as the convex hull of \(\supp f\cup \set{(0,0)}\).
\end{definition}

For the proof we will require some basic knowledge of the theory of derivations, which we introduce here only in the context of polynomial rings over a field.
For a more general treatment see \autocite{vdE2000} or \autocite{Mat70}.

\begin{definition}[Derivation]
  Let \(K\) be a field of characteristic 0, and \(R=K[x_1,\ldots,x_n]\). A \(K\)-derivation on \(R\) is a \(K\)-linear map \(D\colon R\to R\) that satisfies Leibniz's rule, i.e.
  \[D(fg)=fD(g)+D(f)g\qquad\forall f,g\in R.\]
  If for all \(f\in R\) there exists an integer \(n\) such that \(D^n(f)=0\), the derivation is called \emph{locally nilpotent}.
\end{definition}
\begin{example}\label{ex:der}\hfill
  \begin{enumerate}
    \item The partial derivative \(D(f)=f_x\) is a locally nilpotent derivation on \(K[x,y]\). In fact, it is easy to see that the derivations on \(R=K[x_1,\ldots,x_n]\)
          form a free \(R\)-module with basis \(\set{D_i(f)=f_{x_i}}_{i=1}^n\).
    \item For any \(f\in K[x,y]\) the map \(D(h)=J(f,h)\) is a derivation. It is not, however, locally nilpotent in general. Consider the special case, where
          \(R=K[x,y]\) and \(f=x^ay^b\). We now show that if \(D(h)=J(x^ay^b,h)\) is locally nilpotent, then \(ab=0\), as this will be useful to us later.
          Suppose that \(0<a\le b\) and
          observe that
          \[D(x^sy^t)=\det
            \begin{pmatrix}
              ax^{a-1}y^b & bx^ay^{b-1} \\
              sx^{s-1}y^t & tx^sy^{t-1}
            \end{pmatrix}=(at-bs)x^{a+s-1}y^{b+t-1}.
          \]
          This immediately shows that \(D\) cannot be locally nilpotent (e.g. \(D^n(y)\ne 0\) for any \(n\ge 0\)).
  \end{enumerate}
  % \begin{remark}
  %   One can prove that the derivation \(D(h)=J(f,h)\) is locally nilpotent if and only if \(f=P(u)\), where \(u\) comes from an automorphism of \(R=K[x,y]\) and \(P\) is a polynomial.
  %   It can be deduced from a theorem of Rentschler that says that the kernel of a derivation \(R^D\coloneqq \set{f\in R\mid D(f)=0}\) is generated over \(K\) by a single polynomial.
  %   For the proof of Rentschler's theorem see for example \autocite{vdE2000}, it can also be quite easily deduced from \cref{lem:triangle} and \cref{lem:psi}.
  % \end{remark}
\end{example}

Suppose that \((f,g)\) is an automorphism of \(R\) and \((s,t)\) is the inverse automorphism. Then \(J(f,g)J(s,t)=1\), so \(J(f,g)=c\in K^*\). Multiplying \(f\) by \(c^{-1}\) we can always
assume that \(J(f,g)=1\), which we do from now on. Two polynomials satisfying \(J(f,g)=1\) are called Jacobian mates and are of independent interest. We now prove a result about
derivations obtained from a pair of Jacobian mates.
\begin{proposition}\label{prop:jac-mates}
  Suppose that \((f,g)\) is an automorphism of \(R\). Then \(D(h)\coloneqq J(f,h)\) is a locally nilpotent derivation on \(K[f,g]=K[x,y]\).
  Moreover, \(D_w(h)\coloneqq J(f_w,h)\) is a locally nilpotent derivation on \(K[x,y]\).
\end{proposition}
\begin{proof}
  The first statement follows since if \(J(f,g)=1\), then \(D(h)=J(f,h)\) is just the standard partial derivative with respect to the second variable on \(K[f,g]\). For the second statement note that
  for any \(u,v\in R\)
  if \(J(u_w,v_w)\ne 0\), then \(J(u_w,v_w)=J(u,v)_w\). 
  Applying this repeatedly to \(D_w\) we get that 
  \(D_w^k(h_w)=D^k(h)_w\) for any \(k\ge 0\), unless the expression on the left vanishes. Since \(D\) is locally nilpotent 
  this means that \(D_w(h)\) is locally nilpotent on \(K[f,g]_w\coloneqq K[p_w\mid p\in K[f,g]]\). However, since \(K[f,g]=K[x,y]\), we have
  \(K[f,g]_w=K[x,y]_w=K[x,y]\).
\end{proof}

\begin{proposition}\label{prop:hom-poly}
  Let \(p,q\in R\) be homogeneous polynomials such that \(J(p,q)=0\). Then there exists \(h\in R\), \(\alpha,\beta\in K\), and integers \(a,b\ge 0\) such that \(p=\alpha h^a\) and \(q=\beta h^b\).
\end{proposition}
\begin{proof}
  Suppose \(\deg(p)=n\) and \(\deg(q)=m\). Let \(f=p^m\) and \(g=q^n\), so that \(\deg(f)=\deg(g)=nm\). If either \(n\) or \(m\) is zero, the theorem is trivial, so assume \(n,m\ge 1\).
  Notice the following property of homogeneous polynomials,
  sometimes called Euler's formula:
  \[xf_x+yf_y=nmf\quad \text{and}\quad xg_x+yg_y=nmg.\]
  In matrix form this can be written as
  \[
    \begin{pmatrix}
      f_x & f_y \\
      g_x & g_y
    \end{pmatrix}\vec{x}{y}=nm\vec{f}{g}.
  \]
  Since \(J(f,g)=0\) the Jacobian matrix of \(f\) and \(g\) has rank 1. Since the equation above has solutions this means that the augmented matrix
  \[\begin{pmatrix}
      f_x & f_y & nmf \\
      g_x & g_y & nmg
    \end{pmatrix}\]
  also has rank 1, so \(f_xg-fg_x=0=f_yg-fg_y\). From this we conclude that the partial derivatives of the rational function \(f/g\) vanish, so \(f=cg\) for some \(c\in K\).
  This means \(p^m=cq^n\), and the proposition follows by examining the decomposition of \(p\) and \(q\) into irreducible factors.
\end{proof}
Finally, we need to recall one more fact regarding the logarithmic derivatives of polynomials.
\begin{proposition}\label{prop:poles}
  Suppose \(f\in K(x),\, \pi\in K[x]\) with \(\pi\) irreducible, and let \(\nu\) be the \(\pi\)-adic valuation on \(K(x)\). If \(\nu(f)=k\in \Z\), then 
  \[\frac{f'}{f}=k \frac{\pi'}{\pi}+g\quad \text{for some }g\in K(x)\text{ such that }\nu(g)\ge 0.\]
\end{proposition}
\begin{proof}
  Write \(f=\pi^k u\) with \(\nu (u)=0\). Then by direct computation
  \[\frac{f'}{f}=k\frac{\pi'}{\pi}+\frac{u'}{u},\]
  and \(\nu(u'/u)\ge 0\).
\end{proof}

%%%%%%%%%%%%%%%%%%%%%%%%%%%%%%%%%%%%%%%%%%%%%%%%%%%%%%%%%%%%%%%%%%%%%%%%%%%%%%%%%%%%%%%%%%%%%%%%%%%%%%%%%%%
\section{Proof of Jung's theorem}\label{sec:3}
In this section we prove Jung's theorem: every automorphism of \(K[x,y]\) is a composition of linear automorphisms and elementary automorphisms
(given by \(x\mapsto x\), \(y\mapsto y+p(x)\) or \(x\mapsto x+p(y)\), \(y\mapsto y\) for a polynomial \(p\)).
Let \((f,g)\) be an automorphism of \(R\) and let \(m=\deg_x(f)\) and \(n=\deg_y(f)\). Without loss of generality (possibly applying the automorphism \((y,x)\)) we can suppose that \(n\ge m\).
We first prove the following lemma on the shape of the Newton polygon \(\cN(f)\).
\begin{lemma}\label{lem:triangle}
  The Newton polygon \(\cN(f)\) is a triangle with vertices \((m,0),(0,n)\), and \((0,0)\).
\end{lemma}
\begin{proof}
  Suppose there is a vertex \((a,b)\) of \(\cN(f)\) with \(a,b>0\). Without loss of generality we can assume that \(a\le b\). Choose a weight \(w\) such that \(f_w=cx^ay^b\).
  By \cref{prop:jac-mates} we know that the derivation
  \(D(h)=J(x^ay^b,h)\) is a locally nilpotent derivation on \(R\). This, however, leads to a contradiction with \cref{ex:der} (ii).
\end{proof}
This illustrates a clear strategy: compose \((f,g)\) with an elementary automorphism to reduce the \emph{area} of the triangle \(\cN(f)\).
If we can get the area to vanish, it means that
\(m=0\), i.e. \(f=f(y)\). Since then \(J(f,g)=-f_yg_x=1\), this implies \(\deg_y(f)=1\) and \(g=g_0(y)+cx\) and \((f,g)\) clearly is a composition of linear and elementary automorphisms,
finishing the proof. Before proceeding we prove a small corollary that helps prove the theorem if \(K\) is not algebraically closed.
\begin{corollary}\label{cor:non-acf}
  Let \((f,g)\in \Aut R\).
  The leading form with respect to the weight \((1,1)\) of \(f\) is a power of exactly one (linear) prime factor in \(K[x,y]\).
\end{corollary}
\begin{proof}
  If \(f_w\) is a monomial the result holds trivially.
  Suppose this is not the case and let \(f_w=(\alpha x+\beta y)(\gamma x+\delta y)f_1(x,y)\)
  for some \(\alpha,\beta,\gamma,\delta\in \bar{K}\) and \(f_1\in \bar{K}[x,y]\).
  Since the two linear factors are distinct, \((\alpha,\beta)\) and \((\gamma,\delta)\) are linearly independent,
  and by a linear automorphism we can set \(s=\alpha x+\beta y\), \(t=\gamma x +\delta y\).
  We now have \(f_w(s,t)=s^at^bf_1(s,t)\). Since \(a,b\ge 1\), this means that \(\cN(f(s,t))\) is not a triangle. This contradicts \cref{lem:triangle}.
  This immediately implies that for \(f_w=\alpha(\beta x+\gamma y)^k\) the scalars \(\alpha,\beta,\gamma\) lie not only in \(\bar{K}\) but in \(K\). Indeed, let \(L\) be the splitting field
  of a dehomogenized polynomial \(f_w(x,1)\). The Galois group of \(L/K\) acts trivially (as \(f_w(x,1)\) has only one root), so automatically \(L=K\).
\end{proof}

Let \(A(f)=\deg_x(f)\deg_y(f)\). Suppose that \(0<m\le n\). We first dispose of the case \(m=n\). By \cref{cor:non-acf}, the leading form of \(f\) with respect to the total degree is
\[f_{(1,1)}=cL^m\]
for some \(c\in K^*\) and a linear form \(L\in K[x,y]\). A linear automorphism \(\zeta\) can send \(L\) to \(x\). Then the leading form of \(\zeta(f)\) is \(cx^m\), and therefore
\[\deg_x(\zeta(f))=m,\qquad \deg_y(\zeta(f))<m.\]
In particular, \(A(\zeta(f))<A(f)\).

It remains to consider the case \(0<m<n\). Choose a primitive integral weight \(w\) orthogonal to the edge \(e\) of \(\cN(f)\) connecting \((m,0)\) and \((0,n)\). Write
\[\alpha=w(x),\qquad \beta=w(y).\]
Since \(\alpha m=\beta n\) and \(m<n\), we have \(\alpha>\beta>0\). Let \(f_w=c\varphi^d\), where \(c\in K^*\), \(\varphi\in K[x,y]\), and \(d\) is maximal. We will deduce the required degree reduction from the following crucial lemma.
\begin{lemma}\label{lem:psi}
  Suppose that \(w\) is a weight function on \(R\) such that \(w(x),w(y)>0\).
  There exists a \(w\)-homogeneous polynomial \(\psi\) that satisfies
  \begin{equation}
    J(\varphi,\psi)=\varphi.     \tag{\(\star\)}
  \end{equation}
\end{lemma}
\begin{proof}
  By \cref{prop:jac-mates} the derivation \(D(h)\coloneqq J(f_w,h)\) is locally nilpotent on \(R\).
  Note that \(D(x)=-(f_w)_y\ne 0\), so 
  pick \(k\ge 1\) such that \(D^k(x)\ne 0\), but \(D^{k+1}(x)=0\) and let \(\eta=D^{k-1}(x)\). By \cref{prop:hom-poly},
  since \(\varphi,D(\eta)\) are homogeneous we conclude that \(D(\eta)=c_1\varphi^{d_1}\).
  Therefore, \(J(f_w,\eta)=J(\varphi^d,\eta)=c_1\varphi^{d_1}\), which means that \(J(\varphi,\eta)=c_2\varphi^{d_1-d+1}\). Now for \(\psi=c_2^{-1}\eta\varphi^{d-d_1}\) we have \(J(\varphi,\psi)=\varphi\).
  A priori we only know that \(\psi\in K(x,y)\), it remains to check that it is in fact a polynomial. Write
  \[\varphi=x^rp(z),\quad \psi=x^sq(z)\qquad \text{for} \ z=x^{-\frac{w(y)}{w(x)}}y,\, r=\frac{w(\varphi)}{w(x)},s=\frac{w(\psi)}{w(x)}.\]
  Now \((\star)\) is equivalent to
  \[w(\psi)p'q-w(\varphi)pq'=cp.\]
  Dividing by \(pq\) we get 
  \[\tag{\(\dagger\)}w(\psi)\frac{p'}{p}-w(\varphi)\frac{q'}{q}=\frac{c}{q}.\]
  Suppose \(\nu(p)=\a\ge 0\) and \(\nu(q)=-\b < 0\).
  Using \cref{prop:poles} we can write 
  \[\frac{p'}{p}=\a \frac{\pi'}{\pi}+p_1\quad\text{and}\quad \frac{q'}{q}=-\b \frac{\pi'}{\pi}+q_1\]
  for some \(p_1,q_1\in K(z)\) with \(\nu(p_1),\nu(q_1)\ge 0\).
  Combining this with (\(\dagger\)) we get that 
  \[(\a w(\psi)+\b w(\varphi))\frac{\pi'}{\pi}+w(\psi)p_1-w(\varphi)q_1=\frac{c}{q}.\]
  The only way for the valuations of both sides to be equal is if the coefficient in front of \(\pi'/\pi\)  vanishes. This is impossible since all of \(\b,w(\varphi),w(\psi)\) are strictly positive, 
  and \(\a\ge 0\).
  We conclude that \(q\) must be a polynomial in \(z\), so in \(y\), hence also \(\psi\) is a polynomial in \(y\). 
  By a similar argument with \(x\) and \(y\) interchanged \(\psi\) is a polynomial in \(x\).
\end{proof}
By \cref{prop:jacob-w}, equation \((\star)\) forces
\[w(\psi)=w(x)+w(y)=\alpha+\beta.\]
Consequently, every monomial \(x^iy^j\) appearing in \(\psi\) satisfies
\begin{equation}\label{eq:psi-support}
  \alpha i+\beta j=\alpha+\beta.
\end{equation}
Because \(\alpha>\beta>0\), the only possible monomial involving both variables is \(xy\), there is no monomial on the \(x\)-axis, and there is at most one monomial on the \(y\)-axis. The latter exists precisely when \(\alpha=k\beta\) for an integer \(k\ge 2\), in which case it is \(y^{k+1}\).

The polynomial \(\psi\) cannot be a monomial. Indeed, a pure power of \(y\) is impossible, since then the left-hand side of \((\star)\) has smaller \(x\)-degree than \(\varphi\). Thus a monomial \(\psi\) would have to be \(cxy\). Equation \((\star)\) would then force every monomial \(x^ry^s\) of \(\varphi\) to satisfy \(c(r-s)=1\). Such monomials have distinct \(w\)-degrees, so the \(w\)-homogeneity of \(\varphi\) would make \(\varphi\) a monomial. This is impossible because \(f_w=c\varphi^d\) contains both endpoint monomials \(x^m\) and \(y^n\). It follows from \eqref{eq:psi-support} that
\[\psi=c_1y(x+\lambda y^k)\]
for some \(c_1,\lambda\in K^*\), where \(\alpha=k\beta\).

Consider the elementary automorphism
\[\zeta=(x-\lambda y^k,y).\]
It preserves the weight \(w\), has Jacobian determinant \(1\), and satisfies \(\zeta(\psi)=c_1xy\). Applying \(\zeta\) to \((\star)\) gives
\[J(\zeta(\varphi),c_1xy)=\zeta(\varphi).\]
Since \(\zeta(\varphi)\) remains \(w\)-homogeneous, the same monomial calculation shows that it is a monomial. Hence
\[\varphi=c_2(x+\lambda y^k)^a y^b\]
for some \(c_2\in K^*\) and integers \(a,b\ge 0\). The endpoint monomial \(x^m\) of \(f_w=c\varphi^d\) forces \(b=0\) and \(ad=m\). The maximality of \(d\) then forces \(a=1\). After absorbing constants, we obtain
\[f_w=c(x+\lambda y^k)^m,\qquad n=km.\]

Finally, \(\zeta(f)_w=\zeta(f_w)=cx^m\). Every monomial \(x^iy^j\) of \(\zeta(f)\) satisfies
\[\alpha i+\beta j\le \alpha m=\beta n.\]
Thus \(\deg_x(\zeta(f))=m\) and \(\deg_y(\zeta(f))<n\), so again
\[A(\zeta(f))<A(f).\]
Before each application of this reduction, we may interchange \(x\) and \(y\) so that \(\deg_x(f)\le \deg_y(f)\); this does not change \(A(f)\). We have therefore found, whenever \(A(f)>0\), a linear or elementary automorphism over \(K\) that strictly decreases \(A(f)\). Induction on this nonnegative integer produces an automorphism \(\xi\), composed of linear and elementary automorphisms, such that \(\deg_x(\xi(f))=0\). As explained above, this finishes the proof.

\printbibliography
\end{document}
